\cotefile{157-2}
\scanpage{368}
\mspage{107}
% archive sheet number: 367
$C_\infty$ et $H_C$ sont ``bons'', \unsure{il|ils} nous \unsure{permettront|permettent} de décrire un \unsure{foncteur|facteur} important qui domine la situation envisagée. Moyennant une hypothèse très forte, mais \deleted{\unsure{souvent}} satisfaite dans des cas importants, on peut dans ces deux formules remplacer $\mathrm{Cat\,hot}$ par $\mathrm{Hot}$.

\deleted{Par contre, si} Par contre, sauf pour le cas où $S$ est réduit à un point et où $\mathrm{Hom}_{\mathrm{Cat\,hot}}$ ou $\mathrm{Hom}_{\mathrm{Hot}}$ de $S$ dans $C$ n'est autre que $\pi_0(C)$ $(\simeq\pi_0(C_\infty)\simeq \pi_0(H_C))$, donc effectivement ``bon'', je ne sais pas si
\[
\underset{\substack{\Vert\\\pi_0\underline{\mathrm{Hom}}(S,C)}}{\mathrm{Hom}_{\mathrm{Cat\,hot}}(S,C)}
\quad\text{ou}\quad\mathrm{Hom}_{\mathrm{Hot}}(S,C)
\]
est lié de façon raisonnable aux \deleted{\unsure{deux|deux grands|dans quelque}} ensembles précédents (déterminés par $H_C$, resp.\ \inserted{par} $C_\infty$), \deleted{et} \deleted{s'ils sont} p. ex. s'ils lui sont isomorphes de façon naturelle, et s'ils ont une signification géométrico-topologique raisonnable. Je commence à soupçonner que \uline{non} -- \unsure{alors qu'on|dans quoi on} aurait ici un cas où une relation \deleted{\unsure{entre}} de type très général concernant un $\pi_0$, ici donc
\[
\pi_0(C)\simeq\pi_0(\unsure{C|L}_\infty)\simeq\pi_0(H_C),
\]
ne se prolonge pas en un iso. sur les $\pi_i$ supérieurs, et en une \deleted{\unsure{bonne}} propriété de \unsure{hoto-équivalence|hoto-équivalences}.
